\subsection{The Second Sunday: \latin{Populus Sion}}\label{sec:second-sunday}

In 2026 the Second Sunday falls on 6 December. The station is at Holy Cross in Jerusalem. The Mass opens
\latin{Populus Sion} (gloss: people of Sion), and its Collect is the second of the four
\latin{Excita} prayers: \latin{Excita, Domine, corda nostra} (gloss: stir up our hearts, O
Lord). The readings are St Paul to the Romans (15:4--13) and the Gospel of John the
Baptist's question from prison (Matthew 11:2--10).

\subsubsection{Hope from the Scriptures}

The Epistle begins with a reason for reading the Old Testament at all:

\begin{quote}
For what things soever were written were written for our learning: that, through patience and
the comfort of the scriptures, we might have hope.\par
{\smallskip\noindent\small ---Romans 15:4, Douay-Rheims (Challoner).\par}
\end{quote}

\noindent Between that verse and the Epistle's last, Paul gathers texts of the Law, the Psalms
and the Prophets in which the Gentiles praise God, the last of them Isaiah's promise of the root
of Jesse, who rises to rule the nations.

The Office of the same Sunday reads that prophecy. The first lessons at Matins are from Isaiah
11, \latin{Et egredietur virga de radice Jesse}:

\begin{quote}
And there shall come forth a rod out of the root of Jesse, and a flower shall rise up out of his
root.

And the spirit of the Lord shall rest upon him: the spirit of wisdom, and of understanding, the
spirit of counsel, and of fortitude, the spirit of knowledge, and of godliness.\par
{\smallskip\noindent\small ---Isaias 11:1--2, Douay-Rheims (Challoner).\par}
\end{quote}

\noindent Before 1960 the Breviary read in the second nocturn St Jerome's commentary on the first
verse:

\begin{quote}
We, however, understand that the rod out of the root of Jesse signifieth the holy Virgin Mary.
She was a clean stem that had as yet put forth no shoot; as we have read above: ``Behold, the
Virgin shall conceive and bear a son.'' (Isa. vii. 14.) And the flower we believe to mean the
Lord our Redeemer, Who hath elsewhere compared Himself to a flower; ``I am a flower of the plain,
and a lily of the valleys.'' (Cant. ii. 1.)\par
{\smallskip\noindent\small ---Jerome, \textit{Commentary on Isaiah} IV (on Isaiah 11:1), tr.
Bute, \textit{The Roman Breviary} (1908), Winter, p.~228.\par}
\end{quote}

The 1962 Office no longer reads Jerome's lesson on this Sunday, but the Isaiah it reads is the
same. In 2026 the Tuesday of this week is the feast of the Immaculate Conception, so the Virgin
whom Jerome saw in the stem is honored two days after the Sunday that reads the prophecy. In a
year when 8 December falls on a Sunday of Advent, the 1960 rubrics give the feast precedence over
the Sunday (n.~15).

The law of the feast is a separate matter from its rank. Under the Code of Canon Law of 1983, the
law of the Latin Church in 2026, the Immaculate Conception is among the holy days that must be
observed besides Sunday (c.~1246 \S\,1), and ``On Sundays and other holy days of obligation, the
faithful are obliged to participate in the Mass'' (c.~1247, in the English published by
the Holy See, an unofficial working aid; the Latin governs). The same Code lets a conference of bishops, with the
prior approval of the Apostolic See, suppress some of the holy days of obligation or transfer them
to a Sunday (c.~1246 \S\,2). Whether 8 December binds in a given country is that conference's
law, which this book does not state.

\subsubsection{The question from prison}

In the Gospel John, in prison, hears of Christ's works and sends two of his disciples with a
question:

\begin{quote}
Art thou he that art to come, or look we for another?\par
{\smallskip\noindent\small ---Matthew 11:3, Douay-Rheims (Challoner).\par}
\end{quote}

\noindent Christ answers
with his works: the blind see, the lame walk, the poor have the gospel preached to them. Then he
adds a warning, and turns to praise John. The 1962 Office reads at Matins the opening of St
Gregory's homily on the passage:

\begin{quote}
The sight of so many signs and so many mighty works should have been a source of wonder, and not
a stumbling-block. And yet the unfaithful found these very works a rock of offence, when they
afterwards saw Him Who had worked so many miracles dying on the Cross. Hence Paul saith: ``We
preach Christ crucified, unto the Jews a stumbling-block and unto the Gentiles foolishness.'' (1
Cor. i. 23.) It is indeed folly in the eyes of men to say that the Author of life died for men:
and thus men put as a stumbling-block to hinder them from coming to JESUS, the very thing that
doth oblige them the most unto Him. For the more humbling God hath undergone for man's sake, the
more worthy is He that man should worship Him.\par
{\smallskip\noindent\small ---Gregory the Great, \textit{Homiliae in Evangelia} 6.1, tr. Bute,
\textit{The Roman Breviary} (1908), Winter, p.~229. Bute heads the homily as the tenth; it is
the sixth in the standard numbering.\par}
\end{quote}

The homily goes on to the warning itself:

\begin{quote}
And blessed is he that shall not be scandalized in me.\par
{\smallskip\noindent\small ---Matthew 11:6, Douay-Rheims (Challoner).\par}
\end{quote}

\noindent Gregory hears in it the Passion foretold:

\begin{quote}
Now what is this, but a plain mention of that time, when He afterwards humbled Himself, becoming
obedient unto death, even the death of the Cross? It is as if He said: ``I indeed do wonderful
works, but the day will come when I shall not refuse to suffer shame and evil treatment. Take
heed then, ye who now worship Me for the works' sake, that when I come to die ye despise Me not
for My death's sake.''\par
{\smallskip\noindent\small ---Gregory the Great, \textit{Homiliae in Evangelia} 6.1, tr. Bute,
\textit{The Roman Breviary} (1908), Winter, p.~229.\par}
\end{quote}

Read in Advent, the warning reaches back to the manger: the Child awaited at Christmas is the one
who will die on the Cross, and the humility of the one is of a piece with the humility of the
other.

When John's disciples have gone, Christ asks the crowds what they went out into the desert to
see: a reed shaken by the wind? Gregory, as the \textit{Catena Aurea} gives his homily, takes
the reed for the carnal mind, which a breath of praise bends one way and a breath of detraction
the other. John was not such a reed: whatever befell him, he stood straight. The figure gives
the week an examination of conscience, how far one's conduct moves with what others say of it.

\medskip
\noindent\textit{Counsel for the week.} Read Romans 15:4--13 on the Sunday and ask, at the end of
it, for the joy and peace in believing that Paul asks for the Romans. Read Isaiah 11:1--10 once
this week, slowly, with Jerome's lesson beside it. On the Tuesday keep the Immaculate
Conception, at Mass if you can whether or not your country binds you to it. Once this week,
before the crucifix, read again the last sentence of Gregory's first passage above and give
thanks for what it says. Watch one day
this week for the moments when praise or blame bends you, and bring what you find to the night
examen (Section~\ref{sec:prayer}).
